\chapter{Asymptotic Structure of Quantum Channels: Supporting Results}
\label{ch:channel_asymptotics_supporting_results}

This appendix supports Chapter~\ref{ch:channel_asymptotics}. It proves the
positivity, trace-preservation, unitality, and trace-adjoint properties of the
mean-ergodic projections used there.

\section{Preservation under the mean-ergodic projection}

This section proves the positivity, trace-preservation, and unitality assertions
used in Theorem~\ref{thm:positive_trace_preserving_mean_ergodic_projection}.
For a bounded-orbit endomorphism $T$, let $P_T$ denote the mean-ergodic
projection of Definition~\ref{def:mean_ergodic_projection}.

\begin{theorem}[Positive mean-ergodic projection]
    \label{thm:positive_mean_ergodic_projection}
    \lean{IsPositiveMap.meanErgodicProjection_isPositiveMap}
    \leanok
    \uses{def:positive_map, def:mean_ergodic_projection,
        thm:mean_ergodic_bounded_orbits}
    Let $T:\MN{D}\to\MN{D}$ be positive and have bounded orbits. Then its
    mean-ergodic projection $P_T$ is positive.
\end{theorem}

\begin{proof}\leanok
    Every Ces\`aro average is positive. Therefore, for $X\ge0$, closedness of
    the positive cone and convergence in
    (\ref{eq:channel_mean_ergodic_limit}) give $P_T(X)\ge0$.
\end{proof}

\begin{theorem}[Trace-preserving mean-ergodic projection]
    \label{thm:trace_preserving_mean_ergodic_projection}
    \lean{IsTracePreservingMap.meanErgodicProjection_isTracePreservingMap}
    \leanok
    \uses{def:trace_preserving, def:mean_ergodic_projection,
        thm:mean_ergodic_bounded_orbits}
    Let $T:\MN{D}\to\MN{D}$ be trace-preserving and have bounded orbits.
    Then its mean-ergodic projection $P_T$ is trace-preserving.
\end{theorem}

\begin{proof}\leanok
    Every nonempty Ces\`aro average preserves the trace. Convergence of the
    averages and continuity of the trace therefore give
    $\tr(P_T(X))=\tr(X)$.
\end{proof}

\begin{theorem}[Identity under the mean-ergodic projection]
    \label{thm:unital_mean_ergodic_projection}
    \lean{LinearMap.HasBoundedOrbits.meanErgodicProjection_one}
    \leanok
    \uses{def:mean_ergodic_projection, thm:mean_ergodic_bounded_orbits}
    Let $T:\MN{D}\to\MN{D}$ have bounded orbits. If $T(\Id)=\Id$, then
    $P_T(\Id)=\Id$.
\end{theorem}

\begin{proof}\leanok
    The identity is a fixed point of $T$, so the fixed-point characterization
    of the mean-ergodic projection gives the conclusion.
\end{proof}

These three preservation statements give the corresponding properties of the
Ces\`aro projection in
Theorem~\ref{thm:positive_trace_preserving_mean_ergodic_projection}.

\section{Trace adjoints of ergodic projections}

This section proves the two adjoint facts used in
Theorem~\ref{thm:positive_unital_adjoint_mean_ergodic_retraction}. For a linear
map $S$, write $S^*$ for its adjoint with respect to the trace pairing; for a
bounded-orbit map $T$, write $P_T$ for its mean-ergodic projection.

\begin{theorem}[Trace preservation and the trace adjoint]
    \label{thm:trace_preserving_iff_trace_adjoint_unital}
    \lean{isTracePreservingMap_iff_traceAdjointMap_one}
    \leanok
    \uses{def:trace_preserving, def:trace_pairing_adjoint}
    On any finite full matrix algebra, a linear endomorphism $S$ is
    trace-preserving if and only if its trace-pairing adjoint fixes the
    identity:
    \begin{align}
        (\forall X,\ \tr(S(X))=\tr(X))
         & \quad\Longleftrightarrow\quad S^*(\Id)=\Id.
        \label{eq:channel_tp_adjoint_unital}
    \end{align}
\end{theorem}

\begin{proof}\leanok
    \uses{thm:trace_pairing_adjoint_identity, thm:trace_nondeg}
    The trace-pairing identity gives
    $\tr(S^*(\Id)X)=\tr(S(X))$. The equivalence
    (\ref{eq:channel_tp_adjoint_unital}) follows from nondegeneracy of the
    trace pairing.
\end{proof}

\begin{theorem}[Adjoint mean-ergodic projection]
    \label{thm:adjoint_mean_ergodic_projection}
    \lean{LinearMap.HasBoundedOrbits.traceAdjointMap_meanErgodicProjection_apply_eq_self_iff,
        LinearMap.HasBoundedOrbits.traceAdjointMap_meanErgodicProjection_apply_self,
        LinearMap.HasBoundedOrbits.range_traceAdjointMap_meanErgodicProjection}
    \leanok
    \uses{def:mean_ergodic_projection, thm:mean_ergodic_bounded_orbits}
    Let $T:\MN{D}\to\MN{D}$ have bounded orbits, let $P_T$ be its
    mean-ergodic projection, and let $P_T^*$ be the trace-pairing adjoint of
    $P_T$. Then
    \begin{align}
        (P_T^*)^2     & =P_T^*, \notag  \\
        \range(P_T^*) & =\ker(T^*-\Id).
        \label{eq:channel_adjoint_mean_ergodic}
    \end{align}
    and $P_T^*(Y)=Y$ if and only if $T^*(Y)=Y$.
\end{theorem}

\begin{proof}\leanok
    \uses{thm:trace_pairing_adjoint_identity, thm:trace_nondeg,
        thm:mean_ergodic_bounded_orbits}
    Since $P_T^2=P_T$, the trace-pairing identity gives
    $(P_T^*)^2=P_T^*$. If $T^*(Y)=Y$, decompose
    $X-P_T(X)=(T-\Id)(Z)$. Then
    \begin{align}
        \tr(Y(X-P_T(X)))
         & =\tr((T^*(Y)-Y)Z)=0.
        \notag
    \end{align}
    Nondegeneracy of the trace pairing gives $P_T^*(Y)=Y$. Conversely,
    $P_TT=P_T$ implies $T^*P_T^*=P_T^*$, so every fixed point of $P_T^*$ is a
    fixed point of $T^*$:
    \begin{align}
        P_T^*(Y)=Y
         & \quad\Longrightarrow\quad
        T^*(Y)=T^*(P_T^*(Y))=P_T^*(Y)=Y.
        \notag
    \end{align}
    Hence
    \begin{align}
        Y\in\range(P_T^*)
         & \quad\Longleftrightarrow\quad P_T^*(Y)=Y
        \quad\Longleftrightarrow\quad T^*(Y)=Y
        \quad\Longleftrightarrow\quad Y\in\ker(T^*-\Id).
        \notag
    \end{align}
    The first equivalence uses idempotence, and the second is the fixed-point
    equivalence proved above. This proves
    (\ref{eq:channel_adjoint_mean_ergodic}).
\end{proof}

The first theorem gives unitality of $P_T^*$ from trace preservation of $P_T$;
the second identifies its range and fixed points, completing the adjoint
retraction argument in
Theorem~\ref{thm:positive_unital_adjoint_mean_ergodic_retraction}.
